\documentclass[11pt]{article}

\usepackage[letterpaper,margin=1in]{geometry}
\usepackage{amsmath,amssymb}
\usepackage{booktabs}
\usepackage{longtable}
\usepackage{array}
\usepackage{xcolor}
\usepackage{listings}
\usepackage{url}
\usepackage{hyperref}
\usepackage{titlesec}
\usepackage{parskip}
\usepackage{fancyhdr}
\usepackage{enumitem}

% Hyperref colors
\hypersetup{
  colorlinks=true,
  linkcolor=black,
  urlcolor=blue!50!black,
  citecolor=black,
  pdftitle={Precision Controller --- Interface Specification},
  pdfauthor={LonghornSilicon},
  pdfsubject={ISA Specification, Version pc-isa-0.1}
}

% Headings
\titleformat{\section}{\normalfont\Large\bfseries}{\thesection.}{0.6em}{}
\titleformat{\subsection}{\normalfont\large\bfseries}{\thesubsection}{0.6em}{}
\titleformat{\subsubsection}{\normalfont\normalsize\bfseries}{\thesubsubsection}{0.6em}{}

% Headers / footers
\pagestyle{fancy}
\fancyhf{}
\lhead{\small Precision Controller ISA}
\rhead{\small \texttt{pc-isa-0.1}}
\cfoot{\thepage}
\renewcommand{\headrulewidth}{0.4pt}

% Code listings
\definecolor{codebg}{gray}{0.97}
\definecolor{codekey}{rgb}{0.13,0.29,0.53}
\definecolor{codecom}{rgb}{0.25,0.5,0.25}
\lstset{
  basicstyle=\ttfamily\small,
  backgroundcolor=\color{codebg},
  frame=single,
  framerule=0.4pt,
  rulecolor=\color{black!30},
  xleftmargin=0.5em,
  xrightmargin=0.5em,
  aboveskip=0.6em,
  belowskip=0.6em,
  showstringspaces=false,
  breaklines=true,
  keywordstyle=\color{codekey}\bfseries,
  commentstyle=\color{codecom}\itshape,
  numbers=none
}

\setlist{itemsep=1pt,topsep=2pt}

\begin{document}

%==============================================================================
% Title page
%==============================================================================
\thispagestyle{empty}
\begin{center}
\vspace*{1.2in}
{\Huge\bfseries Precision Controller}\\[0.4em]
{\LARGE\bfseries Interface Specification}\\[1.2em]

{\large \textsc{LonghornSilicon LLM Inference Accelerator}}\\[0.4em]
{\large Block 1 of 4 --- Attention Compute Unit (ACU)}\\[2em]

\rule{0.7\textwidth}{0.4pt}\\[1em]

\begin{tabular}{rl}
\textbf{Version}:    & \texttt{pc-isa-0.1} (first public draft) \\
\textbf{Date}:       & 2026-05-13 \\
\textbf{Status}:     & Pre-tape-out stub for compiler integration \\
\textbf{Author}:     & Chaithu Talasila, The University of Texas at Austin \\
\textbf{Reference}:  & \texttt{LonghornSilicon/attention-compute-unit} \\
\end{tabular}\\[1em]
\rule{0.7\textwidth}{0.4pt}

\vfill
\begin{minipage}{0.85\textwidth}
\small
\textbf{Scope.} This document describes the externally-visible interface to the
\texttt{precision\_controller} block as it will appear on the chip, on an
FPGA prototype (Xilinx ZCU102/104), or in the bit-accurate software model.
A compiler backend targeting the precision controller writes against this
interface; the hardware implementation must conform to it. This specification
is stable for the precision controller block (ACU) only and will be unified
with the rest of the LonghornSilicon ISA when the KV Cache Engine, Token
Importance Unit, and Memory Hierarchy Controller blocks land.
\end{minipage}

\vspace{0.3in}
\end{center}

\clearpage

%==============================================================================
% TOC
%==============================================================================
\tableofcontents
\clearpage

%==============================================================================
% 1. Block overview
%==============================================================================
\section{Block Overview}
\label{sec:overview}

The precision controller is a streaming block that consumes one signed
attention score per cycle, asserts an end-of-tile signal on the last score,
and produces a single-bit precision decision (INT8 or FP16) one cycle later.
The decision is computed by the integer rearrangement of an entropy-equivalent
ratio gate, requiring no division and no transcendentals:
\begin{equation*}
\mathrm{decision} \;=\; \bigl(\max(|s|) \ll \log_2 N\bigr) \;>\;
\bigl((\Sigma(|s|) \ll 3) + (\Sigma(|s|) \ll 1)\bigr)
\end{equation*}
or equivalently,
\begin{equation*}
\mathrm{decision} \;=\; \max(|s|) \cdot N \;>\; \Sigma(|s|) \cdot 10
\end{equation*}
where $\mathrm{decision} = 1$ means this tile needs FP16 and
$\mathrm{decision} = 0$ means it can use INT8. The constant $N = \mathrm{BLOCK\_M}
\cdot \mathrm{BLOCK\_N}$ is the synthesis-time tile size; the threshold $T = 10$
is hard-coded in the synthesized netlist via the equivalent shift-and-add
$(\Sigma \ll 3) + (\Sigma \ll 1)$. Changing either requires re-synthesis
(see~\S\ref{sec:syntime}).

\medskip
\noindent\textbf{Latency:} one cycle after \texttt{s\_last} asserts.\\
\textbf{Throughput:} one score per cycle, one decision per $N$ scores.\\
\textbf{Footprint:} 30 flip-flops, $\sim$3.4\,k\,\textmu m\textsuperscript{2} stdcell at
SkyWater Sky130A; projects to $\sim$150\,\textmu m\textsuperscript{2} at TSMC 16FFC.\\
\textbf{Bit-exact reference:} \path{sw/reference_model/precision_controller_ref.py}
(143/143 RTL replay tiles match).

%==============================================================================
% 2. Address space and memory map
%==============================================================================
\section{Address Space and Memory Map}
\label{sec:memmap}

The block exposes a 256-byte AXI-Lite slave register window. All registers
are 32-bit, word-aligned. The base address is set at chip integration time
(e.g., \texttt{0x4000\_0000} on the ZCU102 PYNQ overlay).

\begin{table}[h]
\centering
\small
\caption{AXI-Lite register window (offsets relative to base address).}
\label{tab:regs}
\begin{tabular}{@{}l l c l p{6.8cm}@{}}
\toprule
Offset & Name & Access & Reset & Purpose \\
\midrule
\texttt{0x00} & \texttt{CTRL}              & RW    & 0x0   & bit[0]: \texttt{soft\_reset} (write-1 to pulse); bit[1]: \texttt{enable} \\
\texttt{0x04} & \texttt{STATUS}            & R     & 0x1   & bit[0]: \texttt{idle}; bit[1]: \texttt{decision\_valid}; bit[2]: \texttt{tile\_in\_progress}; bit[3]: \texttt{fifo\_full} \\
\texttt{0x08} & \texttt{INFO\_BLOCK\_M}    & R     & (syn) & Synthesis-time \texttt{BLOCK\_M} \\
\texttt{0x0C} & \texttt{INFO\_BLOCK\_N}    & R     & (syn) & Synthesis-time \texttt{BLOCK\_N} \\
\texttt{0x10} & \texttt{INFO\_N}           & R     & (syn) & $\texttt{BLOCK\_M} \cdot \texttt{BLOCK\_N}$ \\
\texttt{0x14} & \texttt{INFO\_SCORE\_WIDTH}& R     & (syn) & Bits per signed score on \texttt{s\_axis\_scores} \\
\texttt{0x18} & \texttt{INFO\_THRESHOLD}   & R     & (syn) & Always 10 in \texttt{pc-isa-0.1} \\
\texttt{0x1C} & \texttt{INFO\_VERSION}     & R     & 0x100 & bits[15:8] major, bits[7:0] minor \\
\texttt{0x20} & \texttt{TILE\_COUNT}       & R     & 0x0   & Complete tiles processed since last reset \\
\texttt{0x24} & \texttt{LAST\_DECISION}    & R     & 0x0   & bit[0]: most recent decision (1 = FP16) \\
\texttt{0x28} & \texttt{DECISION\_FIFO}    & R     & ---   & Pop one entry; bit[0] = decision; bit[31] = \texttt{valid} \\
\texttt{0x2C} & \texttt{DECISION\_FIFO\_LVL}& R    & 0x0   & Number of decisions waiting in the FIFO \\
\texttt{0x30} & \texttt{IRQ\_MASK}         & RW    & 0x0   & bit[0]: enable interrupt on every decision \\
\texttt{0x34} & \texttt{IRQ\_STATUS}       & R/W1C & 0x0   & bit[0]: decision-available IRQ pending; write 1 to clear \\
\bottomrule
\end{tabular}
\end{table}

\medskip
\noindent\textbf{Conventions.} \texttt{RW} = read/write; \texttt{R} = read-only;
\texttt{W1C} = write-1-to-clear; (syn) = value fixed at synthesis time.
Writes to read-only registers are silently dropped. Reading reserved offsets
returns zero.

%==============================================================================
% 3. Streaming data interfaces
%==============================================================================
\section{Streaming Data Interfaces}
\label{sec:stream}

Two AXI-Stream interfaces carry the actual scalar data. The AXI-Lite window in
\S\ref{sec:memmap} is for control only.

\subsection{\texttt{s\_axis\_scores} --- Score Input (Slave)}
\label{sec:s_axis_scores}

\begin{table}[h]
\centering
\small
\begin{tabular}{@{}l c c p{8cm}@{}}
\toprule
Signal      & Width                  & Direction & Purpose \\
\midrule
\texttt{tdata}  & \texttt{SCORE\_WIDTH} & in        & One signed two's-complement attention score \\
\texttt{tlast}  & 1                     & in        & Assert on the final score of each tile \\
\texttt{tvalid} & 1                     & in        & Handshake: data is valid \\
\texttt{tready} & 1                     & out       & Handshake: block is ready to accept \\
\bottomrule
\end{tabular}
\end{table}

\noindent\textbf{Protocol:}
\begin{itemize}
\item A complete tile is exactly $N$ beats. The block does not validate the
  beat count; feeding fewer than $N$ beats and asserting \texttt{tlast}
  mid-tile produces an undefined decision.
\item \texttt{tlast} is asserted on the $N$th beat. The decision is computed
  combinationally on that cycle's \texttt{tdata} (so the running max/sum
  sees the last score), and the decision is latched on the following rising
  edge.
\item Backpressure: when the block is not ready (e.g., the decision FIFO is
  full), \texttt{tready} deasserts. Upstream must hold \texttt{tdata} and
  \texttt{tlast} stable until the handshake completes.
\end{itemize}

\subsection{\texttt{m\_axis\_decisions} --- Decision Output (Master)}
\label{sec:m_axis_decisions}

\begin{table}[h]
\centering
\small
\begin{tabular}{@{}l c c p{8cm}@{}}
\toprule
Signal      & Width & Direction & Purpose \\
\midrule
\texttt{tdata}  & 1 & out & bit[0]: decision (1 = FP16, 0 = INT8) \\
\texttt{tvalid} & 1 & out & Handshake: decision is valid \\
\texttt{tready} & 1 & in  & Handshake: downstream is ready to accept \\
\texttt{tlast}  & 1 & out & Always 1 (each decision is a one-beat packet) \\
\bottomrule
\end{tabular}
\end{table}

\noindent\textbf{Protocol.} One beat per processed tile. If the downstream
consumer is not ready, the decision queues in the internal FIFO (default depth
16). When the FIFO fills, the block stops accepting scores via \texttt{tready}
on \texttt{s\_axis\_scores}.

%==============================================================================
% 4. Logical operations (compiler-facing)
%==============================================================================
\section{Logical Operations}
\label{sec:ops}

Below are the abstract operations a compiler emits. Each maps to a short
sequence of register writes and stream beats; the C API in \S\ref{sec:capi}
gives the concrete library calls.

\begin{table}[h]
\centering
\small
\begin{tabular}{@{}l l l p{6cm}@{}}
\toprule
Operation                & Inputs         & Outputs        & Description \\
\midrule
\texttt{PC\_QUERY}       & ---            & INFO struct    & Read all \texttt{INFO\_*} registers in one transaction \\
\texttt{PC\_RESET}       & ---            & ---            & Soft reset accumulators + clear decision FIFO \\
\texttt{PC\_ENABLE}      & ---            & ---            & Set bit[1] of \texttt{CTRL} \\
\texttt{PC\_PUSH\_TILE}  & $N$ scores     & ---            & Stream a tile via \texttt{s\_axis\_scores} with \texttt{tlast} on the final beat \\
\texttt{PC\_READ}        & ---            & decision bit   & Pop one decision from the FIFO; blocks if empty \\
\texttt{PC\_READ\_BATCH} & count          & count bits     & Pop up to count decisions; non-blocking \\
\texttt{PC\_STATUS}      & ---            & STATUS bits    & Read the \texttt{STATUS} register \\
\bottomrule
\end{tabular}
\end{table}

\noindent
The block has no other externally-observable side effects. There is no
clock-gating register, no power state machine, and no DFT scan-chain control
at this stub level (those land on the chip-level ISA when the full chip is
defined).

\subsection{Worked Lowering Example}
\label{sec:lowering}

This section walks through how a representative compiler IR operation lowers
to a sequence of ISA operations. The example uses a generic high-level
\texttt{flash\_attention} op that any silicon-agnostic compiler is likely to
have; the binding to a specific compiler IR (MLIR, TVM Relax, ONNX, etc.) is
discussed in \S\ref{sec:bindings}.

Suppose the compiler's IR contains:
\begin{lstlisting}[basicstyle=\ttfamily\small\itshape]
%out = flash_attention(%q, %k, %v)
    : tensor<S x D x f16>, tensor<S x D x f16>, tensor<S x D x f16>
    -> tensor<S x D x f16>
\end{lstlisting}

The lowering proceeds in three stages. \textbf{Stage~1: setup.} The compiler
queries the chip's tile dimensions once per compilation unit:

\begin{lstlisting}[basicstyle=\ttfamily\small]
info = PC_QUERY()                 // -> block_m=64, block_n=64, n=4096
PC_RESET()
PC_ENABLE()
\end{lstlisting}

\textbf{Stage~2: per-tile loop.} For each (BLOCK\_M, BLOCK\_N) sub-block of
the $Q \!\cdot\! K^{\top}$ matrix, the compiler emits a precision-routing
decision before dispatching the actual matmul:

\begin{lstlisting}[basicstyle=\ttfamily\small]
for (tile_i, tile_j) in tiles_of(Q, K):
    scores = Q[tile_i] @ K[tile_j].T         // FP16, computed in shared SRAM
    PC_PUSH_TILE(scores.flatten())           // stream N scores in
    decision = PC_READ()                     // 1 cycle after tlast

    if decision == FP16:
        out_tile = fp16_matmul(scores, V[tile_j])
    else:
        scores_int8 = quantize_int8(scores)
        out_tile = int8_matmul(scores_int8, V[tile_j])
    accumulate(out, out_tile)
\end{lstlisting}

\textbf{Stage~3: batching (optional).} For higher throughput the compiler
can pre-queue many tiles into \texttt{PC\_PUSH\_TILE} and pop decisions
in batches once enough have completed:

\begin{lstlisting}[basicstyle=\ttfamily\small]
for tile in tiles[:k]:
    PC_PUSH_TILE(tile_scores(tile))      // tiles fan into the FIFO
decisions = PC_READ_BATCH(k)              // pop k results at once
dispatch_int8_tiles ([t for t,d in zip(tiles,decisions) if not d], V)
dispatch_fp16_tiles ([t for t,d in zip(tiles,decisions) if     d], V)
\end{lstlisting}

\noindent
The same three stages work whether the compiler's target is the Python
reference model (phase~0), the FPGA prototype (phase~1), or the silicon
chip (phase~3). Only the runtime implementation of \texttt{PC\_*} changes
under the hood --- Python function call, AXI register write, or PCIe MMIO,
respectively. The lowered IR is identical.

\subsection{Compiler Binding Patterns}
\label{sec:bindings}

The five logical operations in Table~\ref{tab:regs} cover every interaction
a compiler needs with the precision controller. How those operations
surface inside a given compiler depends on the compiler's IR. Sketches of
the four most common patterns:

\noindent\textbf{MLIR dialect.} Define a \texttt{lhsi.pc.*} dialect with
five ops corresponding to the table in \S\ref{sec:ops}. The lowering pass
matches a \texttt{flash\_attention} (or vendor-equivalent) op and rewrites
it into the \texttt{lhsi.pc} sequence above plus the appropriate
\texttt{linalg.matmul} variants. Conversion to LLVM IR happens in a
subsequent pass via a runtime call interface.

\noindent\textbf{TVM Relax / TIR.} Register a backend target named
\texttt{"lhsi-pc"}, add a pattern-matcher in the operator-fusion pass
that recognizes attention sub-graphs, and emit calls to TVM-runtime
extern functions named for each \texttt{PC\_*} operation. The TVM
runtime then dispatches to either the Python reference (test mode) or
the C driver (deployed mode).

\noindent\textbf{ONNX --- graph rewriter.} Add an ONNX-Runtime
\texttt{CustomOpDomain} with one op per \texttt{PC\_*}, then write a
graph-transform pass that recognizes the attention pattern and inserts
the custom ops in the right places. The runtime implementation of each
custom op binds to either the Python reference or the C driver.

\noindent\textbf{Custom IR.} If the compiler has its own IR, the
integration point is wherever its codegen emits the lowest-level
hardware-target operations. Replace the call-site for the existing
attention codegen with a sequence of calls to \texttt{PC\_*} plus the
precision-routed matmul kernels.

In every case, the contract on our side is the same: a compiler binding
that emits the ISA operations defined in this document is correct,
testable today against the Python reference, and forward-compatible
through every integration phase.

%==============================================================================
% 5. C API stub
%==============================================================================
\section{C API Stub}
\label{sec:capi}

What a runtime built on top of this ISA looks like. Concrete header sketch:

\begin{lstlisting}[language=C]
/* lhsi_precision_controller.h --- LonghornSilicon ACU driver API. */

#include <stdint.h>
#include <stdbool.h>
#include <stddef.h>

typedef struct {
    uint32_t block_m;
    uint32_t block_n;
    uint32_t n;            /* = block_m * block_n */
    uint32_t score_width;  /* bits per score */
    uint32_t threshold;
    uint16_t isa_major;
    uint16_t isa_minor;
} lhsi_pc_info_t;

typedef struct lhsi_pc_handle lhsi_pc_handle_t;

/* Open / close */
int  lhsi_pc_open (const char *device, lhsi_pc_handle_t **out);
void lhsi_pc_close(lhsi_pc_handle_t *h);

/* Configuration queries (read INFO_* registers) */
int  lhsi_pc_query(lhsi_pc_handle_t *h, lhsi_pc_info_t *out);

/* Control */
int  lhsi_pc_reset (lhsi_pc_handle_t *h);
int  lhsi_pc_enable(lhsi_pc_handle_t *h, bool enable);

/* Streaming --- push a complete tile (N scores, score_width bits each).
   Blocks until the entire tile has been accepted by the block's tready
   handshake. */
int  lhsi_pc_push_tile(lhsi_pc_handle_t *h,
                       const void *scores, size_t stride);

/* Pop a single decision (blocks if FIFO is empty). */
int  lhsi_pc_read_decision (lhsi_pc_handle_t *h, bool *out_fp16);

/* Pop up to `count` decisions; returns number actually popped. */
int  lhsi_pc_read_decisions(lhsi_pc_handle_t *h,
                            bool *out_fp16, size_t count);

/* Diagnostics. */
typedef struct {
    bool idle;
    bool decision_valid;
    bool tile_in_progress;
    bool fifo_full;
    uint32_t fifo_level;
    uint32_t tile_count;
} lhsi_pc_status_t;
int  lhsi_pc_status(lhsi_pc_handle_t *h, lhsi_pc_status_t *out);
\end{lstlisting}

\noindent
Return codes follow the standard Linux convention: \texttt{0} on success,
\texttt{-errno} on failure.

%==============================================================================
% 6. Synthesis-time configuration
%==============================================================================
\section{Synthesis-Time Configuration}
\label{sec:syntime}

The following parameters are baked in at synthesis time and cannot be changed
at runtime. A compiler reads them via the \texttt{INFO\_*} registers and
tailors its code generation accordingly.

\begin{table}[h]
\centering
\small
\begin{tabular}{@{}l c l p{5.5cm}@{}}
\toprule
Parameter         & Default & Range (validated)                & Notes \\
\midrule
\texttt{BLOCK\_M}     & 64  & 16, 32, 64, 128, 256              & Combined product must make $N$ a power of two \\
\texttt{BLOCK\_N}     & 64  & 16, 32, 64, 128, 256              & (same) \\
\texttt{SCORE\_WIDTH} & 8   & 4, 6, 8, 10, 12, 16               & Defines \texttt{tdata} width on \texttt{s\_axis\_scores} \\
\texttt{THRESHOLD}    & 10  & \textbf{10 only} in this version  & The shift-and-add hardware path implements $\times 10$ exactly \\
\texttt{FIFO\_DEPTH}  & 16  & 4, 8, 16, 32                      & Decision output FIFO depth \\
\bottomrule
\end{tabular}
\end{table}

\noindent
A future revision (\texttt{pc-isa-0.2}) will expose \texttt{THRESHOLD} as a
runtime-writable register. The RTL change is straightforward but breaks the
current zero-cost shift-and-add multiplier on the right-hand side of the
inequality.

%==============================================================================
% 7. Bit-Accurate Reference
%==============================================================================
\section{Bit-Accurate Reference}
\label{sec:ref}

The Python model at
\path{sw/reference_model/precision_controller_ref.py} is the canonical
bit-accurate reference for this ISA. A compiler can verify its codegen
against the model directly:

\begin{lstlisting}[language=Python]
from precision_controller_ref import PrecisionController

pc = PrecisionController()
decision = pc.process_tile(scores)   # exactly what the chip will return
\end{lstlisting}

\noindent
The model is verified bit-exact against the RTL testbench's 143 replay tiles
(see \path{sw/reference_model/test_precision_controller_ref.py}). Any
divergence between the model and the chip is a bug in this specification,
not in the chip.

%==============================================================================
% 8. Integration phases
%==============================================================================
\section{Integration Phases}
\label{sec:phases}

This specification supports a staged integration that begins before silicon
exists. The interface described in this document is the stable contract
through every phase.

\begin{table}[ht]
\centering
\small
\renewcommand{\arraystretch}{1.2}
\begin{tabular}{@{}c p{3.0cm} p{4.0cm} p{5.0cm}@{}}
\toprule
Phase & Timeline & Compiler targets & We provide \\
\midrule
0 & now & Bit-accurate Python and C++ reference models & Models + ISA + tests on the \texttt{compiler-integration-prep} branch \\
1 & when ZCU102 / 104 arrives & AXI-Lite + AXI-Stream on Zynq UltraScale+ & Vivado bitstream + PYNQ driver \\
2 & after KV-cache, token-importance, and memory-controller blocks complete & Same AXI interface, more blocks visible & Multi-block FPGA project \\
3 & post-tape-out (2027+) & PCIe-attached custom chip & TSMC 16FFC silicon + Linux driver \\
\bottomrule
\end{tabular}
\end{table}

%==============================================================================
% 9. Open questions
%==============================================================================
\section{Open Questions}
\label{sec:open}

The following are intentionally on the table for collaborators. None of them
are committed in \texttt{pc-isa-0.1}.

\begin{enumerate}
\item Should \texttt{THRESHOLD} be elevated to a runtime register? Cost: one
      register and a parameterized multiplier in place of the current
      shift-and-add. Benefit: per-workload tuning without re-synthesis.
\item Is a single-bit decision stream sufficient, or do we want to also
      emit the raw \texttt{max} and \texttt{sum} values for compiler-side
      calibration?
\item Should the decision FIFO support coalesced reads (e.g., 32 decisions
      packed into one 32-bit word) for high-throughput batching?
\item Are AXI-Lite + AXI-Stream the right interfaces, or should we expose a
      CHI / TileLink / custom protocol that matches the rest of the chip?
\item Should the block expose a small reference-counting feature so the
      compiler runtime can track tile-to-decision mapping when multiple
      attention streams are in flight?
\end{enumerate}

%==============================================================================
% 10. Change log
%==============================================================================
\section{Change Log}
\label{sec:changelog}

\begin{itemize}
\item \texttt{pc-isa-0.1} (2026-05-13): First public draft. Stable for the
      precision controller block only.
\end{itemize}

\vspace{0.5in}
\noindent\rule{\textwidth}{0.4pt}\\
\noindent\small\textit{LonghornSilicon Adaptive Precision Attention --- ACU Interface Specification.
This document is open and may be redistributed. The TSMC 16FFC PDK referenced
in the integration phases is NDA-protected and is not part of this document
or the public repository.}

\end{document}
